\documentclass{article}
\usepackage{graphicx} % Required for inserting images
\usepackage{xcolor}

\usepackage{amsmath}
\usepackage{amsthm}
\usepackage{thmtools} 
\usepackage{cleveref}


% Number equations within sections (instead of chapters)
\numberwithin{equation}{section}

% Declare theorem environments anchored to sections
\declaretheorem[name=Theorem,numberwithin=section]{theorem}
\declaretheorem[name=Lemma,numberlike=theorem]{lemma}
\declaretheorem[name=Proposition,numberlike=theorem]{proposition}
\declaretheorem[name=Corollary,numberlike=theorem]{corollary}
\declaretheorem[name=Definition,numberlike=theorem]{definition}
\declaretheorem[name=Assumption,numberlike=theorem]{assumption}
\declaretheorem[name=Example,numberlike=theorem]{example}
\declaretheorem[name=Remark,numberlike=theorem]{remark}

% For restatable theorems 
\declaretheorem[name=Theorem,numberwithin=section]{restatabletheorem}

\title{Goedel-Prover-V2: Scaling Formal Theorem Proving with Scaffolded Data Synthesis and Self-Correction}
\author{Colomban Duclaux}
\date{\today}

\usepackage[backend=biber,style=alphabetic]{biblatex} % Ensure style matches your preference
\addbibresource{../references.bib} % Link your .bib file

\newcommand{\dupoc}{\texttt{DUPOC}}
\newcommand{\cupod}{\texttt{CUPOD}}
\newcommand{\cimcic}{\texttt{CIMCIC}}
\newcommand{\dimcid}{\texttt{DIMCID}}
\begin{document}

\maketitle
These notes recap the content exposed in \cite{lin2025goedelproverv2scalingformaltheorem}.

\section{Abstract}
They introduce a series of open-source language models that are state-of-the-art in automated theorem proving.
Approach: Expert iteration, RL, curriculum learning, verifier-guided self-correction, model averaging.

\section{Introduction}
Automated theorem proving is a grand challenge for AI, requiring the construction of step-by-step, machine-verifiable proofs in formal languages like Lean. Requires a completely rigorous logical flow.

Most methods today use reasoning models that are huge and cost a lot. Overall, our results demonstrate that the frontier of formal theorem proving can be advanced without
access to extremely large models, vast computational resources, or proprietary technology.

\section{Method}
\subsection{Chain-of-Thought and Self-Correction}
We structure the pipeline so that, after an initial proof attempt, verification
failures are parsed and communicated back into the model as corrective guidance. The model then
generates proof repairs, leading to an iterative self-correction process.

\subsection{Curating formal Statements}
Creation of a dataset of math problems with increasing difficulty (curriculum learning) using LLM. Use of a \textit{formalizer} to translate existing math problems written in natural language into formal Lean statements.

\subsection{Training Algorithm}
\begin{itemize}
    \item Use an expert model to perform large-scale inference and produce an initial SFT dataset $S_1$.
    \item Use $S_1$ to finetune the prover $\to$ SFT-$S_1$.
    \item Use expert model and SFT-$S_1$ to annotate self-correction data $\to S_2$
    \item SFT on $S_2 \to$ SFT-$S_2$.
    \item Model averaging SFT-$S_1$ and SFT-$S_2$ $\to$ SFT-$S_2$-AVG
    \item Scaffolding data synthesis with SFT-$S_2$-AVG to generate $S_3$
    \item SFT of SFT-$S_2$-AVG on $S_3 \to$ SFT-$S_3$
    \item Model averaging $\to$ SFT-$S_3$-AVG
    \item RL on SFT-$S_3$-AVG $+$ model averaging $\to$ final RL-AVG
\end{itemize}

\section{Related works}
\begin{itemize}
    \item \textbf{Formal Theorem Proving as Whole Proof Generation:} Generate entire \textbf{Lean} proofs in a single pass.
    \item \textbf{Formal Theorem Proving using Proof Search:} Incrementally construct proofs by exploring derivations paths with verifier feedback at each step. Requires a lot of computations.
    \item \textbf{Self-Repair and Verifier-Guided Refinement:} Use auxiliary signals to streamline the reasoning process: informal sketches of proof as base, RAG for relevant theorems, self-verification.
\end{itemize}

\section{Summary}
This paper is a nice intro to how an automated prover is designed and trained. Some interesting aspects pop out from the literature, such as \textbf{expert iteration} and \textbf{RL}.
\textbf{Related works} very interesting, could read the paper on \textbf{DeepSeek-Prover}

\printbibliography
\end{document}